% p.13 — الباب السادس في امثال الثلاثي المجرد
\mspage{13}
\bab{الباب السادس في امثال الثلاثي المجرد}
\begin{sarf}{|c|c|c|c|c|c|}
\thead{(1) اول} & \thead{(2) ثانى} & \thead{(3) ثالث} & \thead{(4) رابع} & \thead{(5) خامس} & \thead{(6) سادس}\\\hline
فَتَحَ & ضَرَبَ & نَصَرَ & سَمِعَ & حَسِبَ & كَرُمَ\\
يَفْتَحُ & يَضْرِبُ & يَنْصُرُ & يَسْمَعُ & يَحْسِبُ & يَكْرُمُ\\
فَتْحًا & ضَرْبًا & نَصْرًا & سَمْعًا & حُسْبَانًا & كَرَامَةً\\
فَاتِحٌ & ضَارِبٌ & نَاصِرٌ & سَامِعٌ & حَاسِبٌ & كَرِيْمٌ\\
فُتِحَ & ضُرِبَ & نُصِرَ & سُمِعَ & حُسِبَ & \x\\
يُفْتَحُ & يُضْرَبُ & يُنْصَرُ & يُسْمَعُ & يُحْسَبُ & \x\\
فَتْحًا & ضَرْبًا & نَصْرًا & سَمْعًا & حُسْبَانًا & \x\\
مَفْتُوْحٌ & مَضْرُوْبٌ & مَنْصُوْرٌ & مَسْمُوْعٌ & مَحْسُوْبٌ & \x\\
اِفْتَحْ & اِضْرِبْ & اُنْصُرْ & اِسْمَعْ & اِحْسِبْ & اُكْرُمْ\\
لَا تَفْتَحْ & لَا تَضْرِبْ & لَا تَنْصُرْ & لَا تَسْمَعْ & لَا تَحْسِبْ & لَا تَكْرُمْ\\
مَفْتَحٌ & مَضْرِبٌ & مَنْصَرٌ & مَسْمَعٌ & مَحْسِبٌ & مَكْرَمٌ\\
مِفْتَحٌ & مِضْرَبٌ & مِنْصَرٌ & مِسْمَعٌ & مِحْسَبٌ & مِكْرَمٌ\\
أَفْتَحُ & أَضْرَبُ & أَنْصَرُ & أَسْمَعُ & أَحْسَبُ & أَكْرَمُ\\
فَتْحَةٌ & ضَرْبَةٌ & نَصْرَةٌ & سَمْعَةٌ & حَسْبَةٌ & كَرْمَةٌ\\
فَتَّاحٌ & ضَرَّابٌ & نَصَّارٌ & سَمَّاعٌ & حَسَّابٌ & كَرَّامٌ\\\hline
\thead{تمرين} & \thead{تمرين} & \thead{تمرين} & \thead{تمرين} & \thead{تمرين} & \thead{تمرين}\\\hline
رَفَعَ & قَصَدَ & طَلَبَ & شَرِبَ & نَعِمَ & قَرُبَ\\
طَحَنَ & كَشَفَ & قَتَلَ & لَبِسَ & وَرِعَ & كَثُرَ\\
دَفَعَ & كَسَبَ & نَظَرَ & حَمِدَ & وَرِمَ & بَعُدَ\\
مَنَعَ & حَلَقَ & سَتَرَ & عَلِمَ & يَئِسَ & بَصُرَ\\
\end{sarf}
